\subsection{Integral isotypic projectors with root-dependent denominators}
\label{integral-isotypic-projector-section}

The following arithmetic estimate refines Adriazola's proof of the
isotypic-projector bound in Theorem~12.54 (B7) of the reconstruction.
The use of the extremal projector on $\operatorname{Hom}(W,V)$ and the
partial-trace identity are taken from that argument.
\cite{adriazola-reconstruction}
We retain a separate binomial denominator for each root, retain the
dimension factor in the partial trace, and include the prime~$2$.
The estimate concerns integral representation theory; the arithmetic
theta maps and their coefficient lattices must still be identified
before it can be applied to the global construction.

Let $p$ be any prime, let $K/\mathbb Q_p$ be finite with valuation
ring $\mathcal O$, and normalize $v$ by $v(p)=1$.  Let $V$ be a
finite-dimensional algebraic representation of the split group
$\mathrm{GL}_r$ over $K$.  Call a full $\mathcal O$-lattice
$\Lambda\subset V$ \emph{admissible} if it is the direct sum of
its intersections with the diagonal weight spaces, and every root
divided power $E_{ab}^{(k)}=E_{ab}^k/k!$, for $a\ne b$ and $k\ge0$,
preserves~$\Lambda$.

Fix an occurring highest weight $\lambda$, choose an irreducible
subrepresentation $W\subset V$ of that type, and write
$n_\lambda=\dim_K W$.  Let $e_\lambda$ be the canonical projector
onto the whole $\lambda$-isotypic summand of~$V$.  For the positive
root $\alpha=\epsilon_a-\epsilon_b$, $a<b$, put $h_\alpha=b-a$ and
\begin{align*}
 \mathcal K_{\alpha,\lambda}
   &=\{k\ge1:k\alpha\text{ is a weight of }
                         \operatorname{Hom}_K(W,V)\},\\
 c_{\alpha,\lambda}(p)
   &=\max\left(\{0\}\cup
       \left\{v_p\binom{h_\alpha+k}{k}:
                       k\in\mathcal K_{\alpha,\lambda}\right\}\right),\\
 C_\lambda(p)&=\sum_{\alpha>0}c_{\alpha,\lambda}(p),
 \qquad
 D_\lambda(p)=\max\{0,C_\lambda(p)-v_p(n_\lambda)\}.
\end{align*}
The weight sets are finite and independent of the chosen copy of $W$.
For $r=1$ the sum over roots is zero.

\begin{theorem}[Integral isotypic-projector bound]
\label{integral-isotypic-projector-theorem}
For every admissible lattice $\Lambda$,
\[
             p^{D_\lambda(p)}e_\lambda\Lambda\subseteq\Lambda.
\]
For a subset $\Sigma$ of the occurring irreducible types, the projector
$e_\Sigma=\sum_{\lambda\in\Sigma}e_\lambda$ satisfies
\[
       p^{\max_{\lambda\in\Sigma}D_\lambda(p)}e_\Sigma\Lambda
                  \subseteq\Lambda.
\]
The empty subset has projector zero and denominator exponent zero.
These are sufficient denominator exponents; no optimality is asserted
in general.
\end{theorem}

\begin{proof}
First consider an algebraic representation $U$ with an admissible
lattice $L$, and let $P_U$ be its invariant projector.  It kills all
nonzero weight spaces.  On the zero-weight space $U_0$, the
Asherova--Smirnov--Tolstoy extremal projector is a normally ordered
product over positive roots of the factors
\begin{equation}\label{integral-isotypic-ast-factor}
 P_\alpha\big|_{U_0}
  =\sum_{k\ge0}(-1)^k
       \binom{h_\alpha+k}{k}^{-1}
                  f_\alpha^{(k)}e_\alpha^{(k)}.
\end{equation}
Indeed the usual coefficient of $f_\alpha^k e_\alpha^k$ is
$(-1)^k/[k!(h_\alpha+1)\cdots(h_\alpha+k)]$ on $U_0$.
Replacing both powers by divided powers gives
\eqref{integral-isotypic-ast-factor}.
All denominators here are nonzero positive integers, and the sums
terminate on a finite-dimensional representation.  The product
formula and its annihilation equations are the classical extremal
projector theorem.\cite{ast-extremal-original,tolstoy-extremal-review}

Each root factor preserves $U_0$, and its $k$th term is zero there
unless $k\alpha$ is a weight of $U$.  Define $c_\alpha$ using this
weight set as above.  Divided-power stability shows that the factor
maps $L\cap U_0$ into $p^{-c_\alpha}(L\cap U_0)$, whence their product
has denominator exponent at most $\sum_{\alpha>0}c_\alpha$.
To identify the product with $P_U$ on $U_0$, decompose $U$ into
irreducibles.  Its image is killed by the positive roots and has
weight zero.  A zero-weight highest vector in an irreducible
algebraic representation of $\mathrm{GL}_r$ generates a trivial
representation.  On the trivial summands all root factors are the
identity.  Thus the product is exactly the invariant projector on
$U_0$.  Nontrivial determinant characters have nonzero diagonal
weight and cause no exception.  The integral weight decomposition
therefore gives
\begin{equation}\label{integral-isotypic-invariant-bound}
                p^{\sum_{\alpha>0}c_\alpha}P_U L\subseteq L.
\end{equation}

Put $\Lambda_W=\Lambda\cap W$.  It is a full admissible lattice in
$W$: the intersection is full, and both its weight projections and
its root divided powers are inherited from $\Lambda$ and $W$.
On $H=\operatorname{Hom}_K(W,V)$ put
\[
                 L_H=\operatorname{Hom}_{\mathcal O}
                                      (\Lambda_W,\Lambda).
\]
The integral weight decompositions make $L_H$ weight graded.  If
$X$ is a root operator, left and right composition commute, so
\[
 X^{(k)}\cdot\varphi
    =\sum_{i+j=k}(-1)^j X^{(i)}\circ\varphi\circ X^{(j)}.
\]
Every term preserves $L_H$; hence this lattice is admissible.
Equation~\eqref{integral-isotypic-invariant-bound} gives
$p^{C_\lambda(p)}P_H L_H\subseteq L_H$.

Choose an $\mathcal O$-basis $b_1,\ldots,b_n$ of $\Lambda_W$, with
dual basis $b_1^*,\ldots,b_n^*$, where $n=n_\lambda$.  The exact
partial-trace identity is
\begin{equation}\label{integral-isotypic-partial-trace}
 e_\lambda(v)=n\sum_{i=1}^nP_H(b_i^*\otimes v)(b_i).
\end{equation}
To check it, restrict $v$ to an irreducible summand $U\subset V$.
If $U$ has a different type, $\operatorname{Hom}_K(W,U)$ has no
invariants.  Otherwise choose an intertwiner $\psi:W\to U$.
Under this identification, the invariant projector on
$\operatorname{End}_K(W)$ is $A\mapsto\operatorname{tr}(A)I/n$.
This follows from Schur's lemma and the invariant trace functional.
The right side of \eqref{integral-isotypic-partial-trace} is then
$\sum_i b_i^*(\psi^{-1}v)\psi(b_i)=v$, as required.

For $v\in\Lambda$, every $b_i^*\otimes v$ lies in $L_H$.
The identity and the invariant-projector estimate imply
\[
 e_\lambda\Lambda
       \subseteq n_\lambda p^{-C_\lambda(p)}\Lambda
       =p^{v_p(n_\lambda)-C_\lambda(p)}\Lambda.
\]
This proves the stated bound.  A sum of projectors with integral
denominator exponents at most $D$ also has integral denominator
exponent at most $D$, proving the assertion for~$\Sigma$.
\end{proof}

\begin{corollary}[A bound from the weight width]
\label{integral-isotypic-width-bound}
Suppose each coordinate of all weights of $V$ lies in an interval
of length at most the nonnegative integer $L$.  For $r\ge2$,
\[
 C_\lambda(p)
 \le\sum_{1\le a<b\le r}\max_{0\le k\le L}
                   v_p\binom{b-a+k}{k}
 \le\binom r2\lfloor\log_p(L+r-1)\rfloor.
\]
In particular, $p>L+r-1$ implies that all isotypic projectors
preserve $\Lambda$.  For $r=1$ the denominator cost is zero.
\end{corollary}

\begin{proof}
Every weight of $W$ is a weight of $V$.  Thus $k\alpha$ in
$\operatorname{Hom}(W,V)$ forces $k\le L$.  The factorial valuation
formula writes the binomial valuation as the sum, over $i\ge1$, of
\[
 \left\lfloor\frac{h+k}{p^i}\right\rfloor
 -\left\lfloor\frac h{p^i}\right\rfloor
 -\left\lfloor\frac k{p^i}\right\rfloor.
\]
Each term is zero or one and vanishes for $p^i>h+k$.  Since
$h\le r-1$, this gives the claimed bound.  The zero-root case
does not require evaluating a logarithm.
\end{proof}

For example, take $V=\operatorname{Sym}^a(K^2)\otimes
\operatorname{Sym}^b(K^2)$ and $d=a+b$.  Its constituent
$\det^j\otimes\operatorname{Sym}^{d-2j}(K^2)$,
$0\le j\le\min(a,b)$, has dimension $d-2j+1$.  The positive-root
weights in the relevant Hom representation are $k\alpha$ for
$1\le k\le d-j$.  Consequently the theorem gives the sufficient
denominator exponent
\begin{equation}\label{integral-isotypic-gl2-bound}
 D_j(p)=\max\{0,\lfloor\log_p(d-j+1)\rfloor-v_p(d-2j+1)\}.
\end{equation}
For $a=b=1$ the dimension factor cancels the apparent denominator
at $p=3$, and both projectors are integral there.  At $p=2$ the
bound is one and is sharp on the ordinary tensor lattice, because
the projectors are $(I\pm\mathrm{swap})/2$.

\begin{corollary}[Balanced Fock coefficients at every prime]
\label{integral-isotypic-balanced-fock}
In a polynomial representation built from vector and covector
letters, restricted to any fixed multidegree, put
\[
 v_{\mathrm{BB}}\left(\sum_\beta F_\beta x^\beta\right)
     =\min_\beta\left(v(F_\beta)+\tfrac12v_p(\beta!)\right),
 \qquad \mathcal B=\{F:v_{\mathrm{BB}}(F)\ge0\},
\]
with $v_{\mathrm{BB}}(0)=+\infty$.  The lattice $\mathcal B$ is
admissible for every prime, including $2$.  If $e_\Sigma$ has
exponent $D$ in Theorem~\ref{integral-isotypic-projector-theorem},
then, without any assumption that $\sqrt p\in K$,
\[
       v_{\mathrm{BB}}(e_\Sigma F)\ge v_{\mathrm{BB}}(F)-D.
\]
\end{corollary}

\begin{proof}
The factorial weights make $\mathcal B$ a full weight-graded lattice.
For one vector letter the root action is $x_a\partial_b$.  Write
the two exponents of a monomial as $n,m$.  If $k\le m$, the change
in its coefficient-plus-factorial valuation under the $k$th divided
power is
\[
 \tfrac12v_p\frac{m!}{(m-k)!}
 +\tfrac12v_p\frac{(n+k)!}{n!}-v_p(k!)\ge0,
\]
because both factorial quotients are divisible by $k!$.  If $k>m$
the image is zero.  For a covector letter the action has a minus
sign, which does not affect valuation.  Root actions on separate
letters commute, so their divided powers have the integral
multinomial expansion into products of divided powers.  Thus
$\mathcal B$ is admissible.

For nonzero $F$, extend to $K'=K(\sqrt p)$ and define the balanced
lattice there by the same coefficient inequalities.  Every possible
value of $v_{\mathrm{BB}}(F)$ lies in the value group of $K'$, so
choose $c\in K'$ with $v(c)=-v_{\mathrm{BB}}(F)$.  Apply the
integral theorem over $K'$ to $cF$ and divide the resulting
valuation inequality by this scalar.  Coefficient valuations and
the canonical projector agree with those over $K$, giving the
assertion.  This argument claims no equality between the base
change of $\mathcal B$ and the balanced lattice over $K'$;
such integral lattices need not coincide.
\end{proof}

The root costs and the dimension saving improve the sufficient
denominator bound itself.  They do not provide the source-test
systems, their ranks, the subsequent Fock detection estimates, or
the normalized arithmetic theta maps needed in the global loss
engine.  An independent exact check constructs the projectors as
rational polynomials in the $\mathfrak{sl}_2$ Casimir for
$0\le a\le b\le5$.  The accompanying script verifies 56 projectors
and 280 prime cases, including $p=2$, with 99 strict improvements
over the common coarse bound.  These finite checks illustrate the
theorem and do not replace its proof.
